\begin{table*}[t]
\centering
\caption{Comparative summary of the federated optimization algorithms evaluated in this study.}
\label{tab:alg_comparison}
\footnotesize
\setlength{\tabcolsep}{4pt}
\renewcommand{\arraystretch}{1.15}
\resizebox{\textwidth}{!}{%
\begin{tabular}{p{3.0cm}p{4.0cm}p{4.5cm}p{4.8cm}}
\toprule
\textbf{Aspect} & \textbf{FedAvg} & \textbf{FedProx} & \textbf{SCAFFOLD} \\
\midrule

Core idea &
Standard federated averaging baseline. &
Federated optimization with proximal regularization. &
Federated optimization with drift correction. \\

Main mechanism &
Aggregates local model updates from the selected clients through weighted averaging. &
Adds a proximal term to the local objective to constrain client updates with respect to the global model. &
Uses server-side and client-side control variates to reduce client drift caused by data heterogeneity. \\

Main strength &
Simple, widely used, and communication-efficient. &
Improves local training stability under heterogeneous client data. &
Provides stronger correction for statistical heterogeneity and client drift. \\

Additional state &
No additional correction variables are required. &
No server-client control variates are required, but local objectives are modified. &
Requires maintaining and updating control variates at both server and client levels. \\

Role in this study &
Serves as the reference baseline for analyzing the effect of the join ratio. &
Tests whether proximal regularization improves the performance/cost trade-off. &
Tests whether drift correction leads to more favorable operating points across join ratio values. \\

\bottomrule
\end{tabular}%
}
\end{table*}
